\scanpage{4}
\uline{Lemme 2} Soit \deleted{$\unsure{B}$} $A$ comme dessus, $G$ un $A$-groupe affine, d'algèbre $B$, satisfaisant les conditions ci-dessus. Alors pour \unsure{tt} sous-module de type fini $M$ de $B$, $\exists$ sous-module de type fini $N$ de $B$ tel que $N\supset M$ et tel que

a) $\tilde N=N$ i.e. $N$ saturé --- ce qui implique $N\otimes_A N\hookrightarrow B\otimes_A B$

b) $\Delta N\subset N\otimes N$

\uline{Dém.} Faisons opérer $G\times G$ \deleted{$\unsure{\longrightarrow}$} sur $B$ par
\[
(g,h)f(s)=f(g^{-1}sh).
\]
\deleted{Alors \unsure{le}} Un sous-module $N$ de $B$ est invariant par $G\times e$ (resp $e\times G$) \deleted{\unsure{ssi}} ssi on a
\[
\Delta N\subset N\otimes B\qquad\text{resp.}\quad\Delta N\subset B\otimes N.
\]
Notons que \deleted{grâce à la platitude de $B$, cette} chacune de ces conditions est stable par passage de $N$ à $\tilde N$. (faire l'argument).

Donc, la condition d'être invariant par $G\times G$, qui est la conjonction des deux, est aussi stable par passage de $N$ à $\tilde N$. Or, on sait que $M\subset N$ \inserted{de t.f.} stable par $G\times G$. Par hyp., $\tilde N$ est encore de t.f., \deleted{\unsure{or}} et est encore stable. \inserted{on peut supposer $N=\tilde N$}

Mais \unsure{cela impliq.}
\[
\Delta\tilde N\subset(\tilde N\otimes B)\cap(B\otimes N)=N\otimes N
\]
cqfd.
\mspage{3}
